\begin{hbwarning}{def:p06-relativization-index-erratum}{Index correction in the relativization lemma}
\emph{Editorial correction to the preceding exact statement.}
The hypotheses list \(v_1,\ldots,v_m\) and
\(\varphi_1,\ldots,\varphi_k\).  Its displayed conclusion is therefore read
with
\[
  \bigwedge_{1\leq j\leq m}v_j\in L
  \qquad\text{and}\qquad
  \bigwedge_{1\leq i\leq k}\varphi_i^L,
\]
not with the printed bounds \(j\leq m\) and \(i\leq k\), which would introduce
undeclared \(v_0\) and \(\varphi_0\).  When \(m=0\) or \(k=0\), the
corresponding conjunction remains absent exactly as the source says.  The
following exact proof already uses only the declared one-based lists.
\end{hbwarning}

\noindent\textit{Base-case reminder.}  In the following exact proof, the
nonemptiness witness for \(L\) is supplied by
Definition~\ref{def:p06-def-empty-base}:
\(\varnothing\in L_1=\{\varnothing\}\).  The printed phrase “whatever the
satisfaction clause returns over the empty domain” is superseded by that
direct calculation and does not extend the handbook's satisfaction relation
to empty structures.
